\section{Modelagem Estocástica do PLD}
\label{sec:rl_pld}

A literatura identifica um conjunto de propriedades do PLD que tornam inadequada sua representação por um único valor esperado: sazonalidade marcada pelo ciclo hidrológico, dependência temporal entre períodos consecutivos, distribuição assimétrica com caudas pesadas e limites regulatórios de piso e teto \cite{guerreiro2017,ccee2024a,mme2002}. \citeonline{guerreiro2017} documenta que a assimetria positiva da distribuição, com picos pronunciados nos períodos secos, leva modelos baseados em valor esperado a subestimar sistematicamente o custo de posições deficitárias, reforçando a necessidade de modelos estocásticos com sazonalidade representada de forma estrutural.

\citeonline{reis2013} fornece a referência metodológica central para essa classe de modelos: os Modelos Autorregressivos Periódicos PAR($p$), que atribuem um conjunto distinto de parâmetros autorregressivos para cada mês do ano, acomodando a estacionariedade periódica das séries hidrológicas e de energia. O modelo GEVAZP, desenvolvido pelo CEPEL para geração de séries sintéticas de energia natural afluente no planejamento do SIN, baseia-se nessa mesma classe. \cite{detzel2011} validam um modelo AR(1) parcimonioso para afluências em usinas da região sul do Brasil, confirmando que modelos autorregressivos de baixa ordem capturam a estrutura de dependência temporal relevante sem requerer estimação de muitos parâmetros.

A utilidade desses modelos autorregressivos está, precisamente, em sua capacidade de gerar cenários que alimentem os métodos de otimização. Para o DEQ, os cenários formam a árvore explícita de realizações sobre a qual o problema é resolvido \cite{birge2011}. Para a PDDE, os cenários são amostrados nas trajetórias da etapa \emph{forward} do algoritmo, sem construção explícita da árvore completa \cite{pereira1991,shapiro2009}. Em ambos os casos, a adequação estatística dos cenários à distribuição real do PLD condiciona diretamente a qualidade da política obtida.

Considerando esses requisitos de adequação estatística, o arcabouço PAR(1) é compatível com o problema de portfólio tratado neste trabalho: captura a sazonalidade mensal e a dependência temporal com apenas doze coeficientes autorregressivos, pode ser estimado a partir de séries históricas da CCEE, e alinha-se com a tradição do GEVAZP validada no contexto brasileiro \cite{reis2013,detzel2011}. A estimação em escala logarítmica garante positividade das séries sintéticas, compatível com a natureza do PLD como preço com piso regulatório positivo.
